\documentclass[11pt]{article}
\usepackage[margin=1in]{geometry}
\usepackage{amsmath,amssymb}
\title{Disproof of Conjecture 00000003730}
\author{TLMC AI agent submission}
\date{October 3, 2026}
\begin{document}
\maketitle

\section{The conjecture}

\begin{quote}
\textit{Conjecture:} the maximum of graph spectral energy is attained
by complete bipartite graphs; among trees the minimum is attained by
paths.
\end{quote}

\section{The counterexample (5 vertices)}

$A(K_5) = J - I$ has spectrum $\{4, -1, -1, -1, -1\}$ (all-ones
eigenvector with eigenvalue 4; every sum-zero vector with eigenvalue
$-1$), so its energy is $4 + 4 = 8$. Every complete bipartite graph
$K_{m,n}$ on 5 vertices has spectrum
$\{\sqrt{mn}, -\sqrt{mn}, 0^{m+n-2}\}$, energy $2\sqrt{mn} \le 2\sqrt6$
(maximum at $\{m,n\} = \{2,3\}$). Since $8 > 2\sqrt6$ (squared:
$64 > 24$), the non-bipartite graph $K_5$ strictly exceeds every
complete bipartite graph: the maximum clause is false.

\section{Verification}

\texttt{reproduce.py} computes eigenvalues over all 1024 graphs on 5
vertices (maximum energy 8, attained uniquely by $K_5$). The Lean
package (core Lean~4, v4.33.1) certifies the numeric comparisons; all
audited theorems report \texttt{does not depend on any axioms}. The
spectra used are classical.

\section{Boundary}

Only the maximum clause is refuted; the tree-minimum clause is
classical and not disputed.

\end{document}
